\documentclass{article}

% ICML-compatible draft. The official ICML workshop may require its own wrapper
% around icml2026.sty; replace only the style package if the workshop provides one.
\usepackage{microtype}
\usepackage{graphicx}
\usepackage{subcaption}
\usepackage{booktabs}
\usepackage{amsmath}
\usepackage{amssymb}
\usepackage{mathtools}
\usepackage{amsthm}
\usepackage{enumitem}
\usepackage{xspace}
\usepackage{listings}
\usepackage{algorithm}
\usepackage{algorithmic}
\usepackage{hyperref}
\usepackage{url}
\usepackage{icml2026}

\icmltitlerunning{RAILS: Resource-Aware Information Lineage for Parallel LLM Agent Search}

\newcommand{\method}{RAILS\xspace}
\newcommand{\keep}{\texttt{KEEP\_EXPERIMENT}\xspace}
\newcommand{\todo}[1]{\textcolor{red!70!black}{[TODO: #1]}}
\newcommand{\tbd}{\textcolor{red!70!black}{TBD}}
\newcommand{\lowerbetter}{$\downarrow$}
\newcommand{\higherbetter}{$\uparrow$}
\newcommand{\eval}{\mathrm{eval}}
\newcommand{\HEAD}{\textsc{head}\xspace}
\newcommand{\head}{\mathrm{HEAD}}
\newcommand{\stale}{\mathrm{stale}}

\lstset{
  basicstyle=\ttfamily\footnotesize,
  breaklines=true,
  frame=single,
  framesep=3pt,
  xleftmargin=0pt,
  xrightmargin=0pt,
  columns=fullflexible,
  keepspaces=true,
  aboveskip=5pt,
  belowskip=5pt
}

\begin{document}

\twocolumn[
\icmltitle{RAILS: Resource-Aware Information Lineage for Parallel LLM Agent Search}

\begin{icmlauthorlist}
\icmlauthor{Anonymous Authors}{anon}
\end{icmlauthorlist}
\icmlaffiliation{anon}{Anonymous Institution}
\icmlcorrespondingauthor{Anonymous Author}{anon.email@domain.com}
\icmlkeywords{LLM agents, black-box optimization, automated research, hyperparameter search, data mixture optimization}
\vskip 0.3in
]

\printAffiliationsAndNotice{}

\begin{abstract}
LLM agents are increasingly used as optimizers: they read a search history, propose a code or configuration change, run a fixed harness, and iterate. Serial agent loops exploit information well because every decision conditions on the latest result and every winning edit can become the next baseline, but they waste parallel accelerators. Naive parallel loops use accelerators efficiently, but many decisions are made from stale history, winners do not compound, and managers often over-sample a single promising axis. We formalize this tension as \emph{information debt} in three channels---state, lineage, and coverage---and propose \method, a resource-aware controller that exposes each channel as an explicit knob. \method combines watermark-gated replanning, a cumulative lineage primitive with optional statistical commit gates, and a rationale/coverage ledger injected into every manager call. In a pilot autoresearch run, \method improves over naive parallel search while preserving high throughput; a Qwen SFT data-mixture task and complete ablations are left as planned workshop experiments.
\end{abstract}

\section{Introduction}
\label{sec:intro}

Many useful machine-learning improvements are obtained through a simple experimental loop. A researcher proposes a change, evaluates it with a fixed harness, reads the metric, and decides what to try next. Recent LLM-agent systems automate the proposer in this loop: an agent edits a file, trains or evaluates for a bounded amount of time, records the result, and repeats. This pattern appears in autonomous ML experimentation such as \texttt{autoresearch} \citep{karpathy_autoresearch}, prompt optimization \citep{yang2023opro,promptagent}, compound-system tuning \citep{textgrad,dspy}, and software-engineering agents evaluated on benchmarks such as SWE-bench \citep{swebench}. We refer to this class as \emph{agent-led search}: the objective is a black-box scalar metric, but the search policy is a language model with access to code, free-form rationales, and accumulated experiment notes.

The deployment setting is usually parallel. A user may have an 8-GPU node or a shared cluster allocation; the evaluation harness is slow; and the natural systems question is how to keep all workers busy. However, the most reliable existing agent-led loops are mostly serial. A serial loop has three attractive properties: the agent sees the freshest outcome before proposing the next experiment, the repository baseline advances immediately when an edit wins, and the agent's mental model of what has been tried is a single coherent trajectory. Its failure mode is obvious: with $K$ available accelerators and a single evaluation in flight, $K-1$ workers are idle.

A naive parallel manager-worker loop makes the opposite trade. The manager dispatches $K$ experiments, workers evaluate them, and the manager dispatches a new batch after receiving results. Throughput is high, but information arrives too late to influence decisions inside the batch. Worse, every worker often forks from the same stale baseline, so parallel winners do not automatically compose. In our pilot runs this is not a small engineering detail: the naive parallel configuration produced many more evaluations per hour than the serial loop, yet its best metric remained within seed noise of the starting baseline. The manager proposed plausible experiments; the loss was in the information plumbing.

This paper argues that the serial-vs-parallel choice is the wrong abstraction. Parallelizing an agent-led loop breaks several information channels, and each channel has a different repair cost. We therefore introduce \method, a controller for \emph{resource-aware agent-led search}. \method represents experiments as edits applied to explicit source revisions, tracks how stale each decision is, and lets a manager selectively promote successful experiments into the shared lineage. It also re-injects the manager's own hypotheses and a programmatic coverage summary into every replanning call, so that high throughput does not erase the search history.

Our contributions are as follows. First, we formalize three kinds of \emph{information debt} in parallel agent-led search: state debt, lineage debt, and coverage debt. Second, we present \method, a minimal set of mechanisms that pays down each debt independently: watermark-gated replanning, cumulative lineage, and a rationale/coverage ledger. Third, we provide a concrete experimental protocol for two workshop-scale domains---small-GPT autoresearch and Qwen supervised-finetuning (SFT) data-mixture search---including metrics that measure not only best score but also throughput and information debt. Our current pilot results support the central claim that lineage and memory matter; the remaining SFT and ablation results are explicitly marked as pending.

\section{Agent-Led Search and Information Debt}
\label{sec:problem}

\paragraph{Agent-led search.}
Let $b_t$ denote the source baseline visible to the proposer at decision time $t$. A candidate is a mutation $u_t$---for example, a constant replacement, a data-mixture vector, or a free-form code patch---which yields a runnable source state $s_t = u_t(b_t)$. A harness evaluates $s_t$ and returns a noisy scalar
\begin{equation}
  y_t = f(s_t) + \epsilon_t, \qquad \epsilon_t \sim \mathcal{D},
\end{equation}
where lower is better in our experiments. The proposer is an LLM policy $\pi_\theta$ that conditions on a history $H_t = \{(b_i,u_i,r_i,y_i)\}_{i<t}$, where $r_i$ is the natural-language rationale recorded when $u_i$ was proposed. The goal is to minimize best-so-far score under a wall-clock budget $W$ and a worker budget $K$.

Serial search sets $K=1$ and usually updates $b_{t+1}$ to the winning source state whenever $y_t$ improves the best known metric. Naive parallel search emits a batch $\{u_{t,j}\}_{j=1}^{K}$ from the same $H_t$ and often from the same baseline $b_t$. The batch yields high utilization but introduces lag between a result and the first decision that can react to it.

\paragraph{Information debt.}
We define information debt as the gap between the information available in an ideal serial decision and the information actually available to a parallel decision. We use three measurable components.

\textbf{State debt} counts how many completed or soon-to-complete results were not visible when a task was dispatched. If task $e$ is dispatched at time $d_e$, let $C(d_e)$ be the completed set at dispatch and $C(t_e)$ the completed set when $e$ finishes. A simple proxy is
\begin{equation}
  D_{\mathrm{state}} = \frac{1}{|E|}\sum_{e\in E} |C(t_e)\setminus C(d_e)|.
\end{equation}
Large state debt means the manager commits many proposals before seeing the consequences of earlier ones.

\textbf{Lineage debt} measures how many evaluations run from a baseline that is no longer the current accepted lineage. If $b_e$ is the baseline used by evaluation $e$ and $\head(t)$ is the latest shared baseline at time $t$, then
\begin{equation}
  D_{\mathrm{lineage}} = \frac{1}{|E|}\sum_{e\in E} \mathbf{1}\{b_e \neq \head(t_e)\}.
\end{equation}
A naive batch loop can have zero explicit commits and still have high \emph{effective} lineage debt because all winners are isolated leaves rather than cumulative improvements.

\textbf{Coverage debt} measures collapse of exploration onto a small set of axes. Let $a(e)$ be the edit axis of evaluation $e$, such as \texttt{lr}, \texttt{warmup}, \texttt{aspect\_ratio}, or \texttt{data\_mix.math}. We report normalized entropy
\begin{equation}
  H_{\mathrm{axis}} = -\frac{1}{\log |\mathcal{A}|}\sum_{a\in\mathcal{A}} p(a)\log p(a),
\end{equation}
where $p(a)$ is the fraction of completed evaluations touching axis $a$. Low entropy is not always bad---exploitation is sometimes correct---but low entropy combined with no lineage progress is a diagnostic of manager fixation.

\begin{table*}[t]
\centering
\small
\begin{tabular}{p{0.16\textwidth}p{0.24\textwidth}p{0.26\textwidth}p{0.26\textwidth}}
\toprule
Debt channel & Serial loop behavior & Naive parallel failure & \method mechanism \\
\midrule
State & Replans after every result & Decisions lag by a batch & Watermark-gated asynchronous replanning \\
Lineage & Winning edits become the next baseline & Parallel winners remain isolated leaves & Cumulative lineage with \keep and rebase validation \\
Coverage & Single coherent notebook of attempts & Manager repeatedly samples one axis & Rationale ledger plus axis coverage manifest \\
\bottomrule
\end{tabular}
\caption{Parallel agent-led search is not one problem but three information-flow problems. \method exposes each as a separate resource knob instead of choosing only between serial and naive batch execution.}
\label{tab:debts}
\end{table*}

Table~\ref{tab:debts} summarizes the channels. The key point is that the repair costs differ. State debt is paid with more LLM calls; lineage debt is paid with more careful source control and possible re-evaluation; coverage debt is paid with prompt/context budget and diversity constraints. A resource-aware system should let the user spend these resources separately.

\section{RAILS}
\label{sec:method}

\method is a manager-worker controller wrapped around a task-specific evaluation harness. The manager is an LLM that emits a structured task graph. Workers execute tasks in isolated sandboxes and write metrics. A runtime maintains a database of worker tokens, a git-backed source lineage, and a compact memory ledger. The design goal is not to replace the LLM manager with a hand-coded optimizer; it is to provide the manager with the missing state that serial execution implicitly supplies.

\subsection{Watermark-Gated Replanning}
\label{sec:watermark}

A batch scheduler calls the manager once per wave. A serial scheduler calls the manager after every evaluation. \method interpolates between these extremes with a queue watermark. Let $N_{\mathrm{ready}}$ be runnable tasks waiting for workers and $N_{\mathrm{run}}$ be running tasks. The manager is woken when
\begin{equation}
  N_{\mathrm{ready}} + N_{\mathrm{run}} < \alpha K
  \label{eq:watermark}
\end{equation}
and a minimum interval has elapsed since the last wake-up. Here $K$ is the number of workers and $\alpha$ is a user-set resource knob. Smaller $\alpha$ means more frequent replanning and lower state debt; larger $\alpha$ means fewer LLM calls and more speculative work. In our current implementation, $\alpha=1.0$ wakes the manager roughly on individual completions when the queue is shallow, while $\alpha=1.5$ behaves like batch planning with a small reserve.

The manager context at wake-up contains all newly completed evaluations since the last call, a best-so-far table, the current lineage, and the ledger described below. The manager then emits a JSON task graph. Most tasks are deterministic parameter patches, avoiding per-worker LLM calls; free-form code edits are supported but treated as the expensive fallback path.

\subsection{Cumulative Lineage}
\label{sec:lineage}

Lineage is the most important difference between agent-led search and ordinary parallel hyperparameter tuning. In serial autoresearch, a winning edit changes the future search space because subsequent experiments are applied on top of it. A naive parallel system that evaluates all leaves from a fixed root cannot discover cumulative improvements unless a separate merge mechanism exists.

\method maintains a shared git lineage and lets the manager explicitly promote an experiment into it. The directive is intentionally simple:
\begin{lstlisting}
<!-- KEEP_EXPERIMENT -->
[{"id":"exp_042_aspect80",
  "reason":"aspect=80 improved bpb and is orthogonal to lr"}]
<!-- /KEEP_EXPERIMENT -->
\end{lstlisting}
On receipt, the runtime looks up the experiment's structured edit, re-applies it to the current shared baseline, checks that the patch still applies, commits the new state, and appends a row to \texttt{lineage.json}. Future evaluations clone from this updated baseline.

\paragraph{Rebase validation for stale winners.}
A parallel winner may have been evaluated on an older baseline. Rather than blindly committing such results, \method distinguishes three states: \emph{validated keep}, \emph{rebase candidate}, and \emph{discard}. If the candidate was evaluated on current \HEAD, it can be committed immediately. If it was evaluated on a stale baseline but its edit applies cleanly to current \HEAD, the runtime can enqueue a short validation run on the rebased source. This converts lineage debt into a bounded validation cost. The current pilot implementation lets the manager decide this manually; the recommended camera-ready version should make rebase validation automatic.

\paragraph{Optional statistical commit gate.}
Because short training runs are noisy, a manager may either keep too aggressively or wait too long. We propose a simple gate for lower-is-better metrics. Let $y_e$ be the candidate score, $y_{\mathrm{HEAD}}$ the score of current \HEAD, and $\hat\sigma$ an estimated run-to-run noise scale from baseline repeats or paired re-evaluations. Commit if
\begin{equation}
  y_{\mathrm{HEAD}}-y_e > \tau_t \hat\sigma, \qquad
  \tau_t = \tau_{\min} + (\tau_{\max}-\tau_{\min})\frac{W-t}{W}.
  \label{eq:commit-gate}
\end{equation}
The threshold is conservative early, when there is time to collect more evidence, and permissive late, when not committing wastes the remaining budget. Equation~\ref{eq:commit-gate} can be used as a tool suggestion to the manager rather than a hard rule, but reporting how often the manager violates the gate gives an interpretable diagnostic.

\subsection{Rationale and Coverage Ledger}
\label{sec:ledger}

Serial agents naturally accumulate a notebook: they remember not only which metric won but why they tried it. A manager-worker runtime must make that notebook explicit. \method stores, for every evaluation, the baseline revision, edit axis, edit value, natural-language rationale, metric, runtime, and failure status. At each manager wake-up it injects a compact ledger consisting of (i) top-performing rows, (ii) most recent rows, (iii) failed or surprising rows, and (iv) the current source files selected by the task plugin.

The ledger has two roles. First, rationales let the manager distinguish hypotheses from accidents. A result such as ``\texttt{aspect\_ratio=80}: $-0.001$'' is less useful than ``\texttt{aspect\_ratio=80}: tests the hypothesis that a slightly wider FFN ratio better matches the hidden dimension; improved by $-0.001$.'' Second, axis summaries make coverage debt visible. The runtime extracts edit axes from structured patches and injects a manifest such as:
\begin{lstlisting}
axis coverage:
- aspect_ratio: tried {64, 80, 96}; best 80
- warmdown_frac: tried {0.55, 0.70}; best 0.55
- final_lr_frac: tried {0.05}; promising, under-sampled
- weight_decay: never tested
\end{lstlisting}
This manifest enables a diversity-constrained manager prompt: each wave should allocate some slots to exploitation of promising axes and some slots to untested or interior-gap axes. In our pilot, rationale logging and source injection are implemented; programmatic coverage manifests are the main planned method upgrade.

\subsection{Task Graph and Execution Modes}
\label{sec:taskgraph}

The manager emits a task graph whose nodes contain an identifier, baseline reference, rationale, resource request, estimated runtime, and an execution mode. In \texttt{param\_patch} mode, the runtime applies structured edits deterministically. Supported edits include module-level constant replacement, regex replacement with an expected match count, anchored block replacement, and unified diff as an escape hatch. In \texttt{code\_edit} mode, a worker LLM performs a free-form code edit inside the sandbox. The latter is useful for architecture changes but expensive and less reproducible.

The task plugin interface consists of runnable source code, a wrapper script that writes \texttt{metrics.json}, a manager prompt, a worker prompt for code-edit fallback, and a workspace generator. This keeps the scheduler and lineage machinery task-agnostic.

\begin{algorithm}[t]
\caption{Resource-aware agent-led search in \method}
\label{alg:rails}
\begin{algorithmic}[1]
\STATE Initialize shared baseline $b_0$, ledger $L$, worker budget $K$
\WHILE{wall-clock budget remains}
  \STATE Launch ready tasks while workers are free
  \STATE Append completed metrics and rationales to $L$
  \IF{$N_{\mathrm{ready}}+N_{\mathrm{run}} < \alpha K$ and wake interval elapsed}
    \STATE Build manager context from current \HEAD, new results, best-so-far table, and coverage ledger
    \STATE Manager emits task graph and optional \keep directives
    \FOR{each \keep directive}
      \IF{candidate evaluated on current \HEAD}
        \STATE Apply edit to shared repo and commit
      \ELSE
        \STATE Rebase edit onto \HEAD and enqueue validation
      \ENDIF
    \ENDFOR
    \STATE Add admitted tasks to ready queue
  \ENDIF
\ENDWHILE
\RETURN best validated source state and full event ledger
\end{algorithmic}
\end{algorithm}

\section{Experiments}
\label{sec:experiments}

The experimental goal is to compare policies at matched wall-clock budgets, not matched numbers of evaluations. We therefore report both optimization quality and resource use. All experiments use the same hardware type within a task, a persistent compile cache when applicable, fixed evaluation wrappers, and fresh source clones per policy.

\subsection{Policies}

\textbf{Serial.} One agent operates in one repository and runs one evaluation at a time. It sees the full history before every proposal and keeps or reverts edits using git. This is the strongest information baseline and the weakest throughput baseline.

\textbf{Naive parallel.} A manager dispatches tasks to $K=8$ workers, but cumulative lineage is disabled. All evaluations fork from a fixed starting baseline unless the task itself changes it. The manager sees summarized results between waves.

\textbf{\method.} The full resource-aware controller with $K=8$, watermark-gated replanning, lineage enabled, source injection, rationale ledger, and deterministic parameter patches. The pilot uses manager-issued keeps; the planned version adds automatic rebase validation and coverage manifests.

\textbf{Ablations.} To identify which channel matters, we ablate cumulative lineage, source injection, rationale ledger, watermark frequency, and coverage steering. Because the SFT task is more expensive, we plan a full ablation only on the autoresearch task and a reduced ablation on SFT.

\subsection{Tasks}

\textbf{Autoresearch hyperparameter/code search.} The task is based on Karpathy's \texttt{autoresearch}, a small single-GPU language-model training setup where the agent edits \texttt{train.py} and optimizes validation bits per byte (\texttt{val\_bpb}) after a fixed five-minute training budget. The search space is intentionally broad: architecture constants, optimizer settings, schedule parameters, batch sizes, and local code changes.

\textbf{Qwen SFT data-mixture search.} The second task fine-tunes Qwen3-0.6B \citep{qwen3} on a weighted mixture of T\"ulu-3 instruction-tuning data \citep{tulu3} grouped into five domains: math, code, chat, instruction following, and reasoning. The manager edits a constrained \texttt{DATA\_MIX} representation while the model, optimizer, and training budget are held fixed. Each evaluation does ten minutes of SFT and reports balanced held-out validation loss across the five domains; secondary metrics include per-domain validation loss and instruction-following evaluation when time allows. Tülu-style open post-training recipes motivate this setup because they make data sources and mixtures explicit.

\subsection{Metrics}

We report five classes of metrics. \emph{Best score} is the best validated metric by the wall-clock deadline. \emph{Best-so-far AUC} integrates the best-so-far curve over time and rewards methods that find improvements early. \emph{Throughput} is completed evaluations per hour and worker utilization. \emph{LLM overhead} is manager calls and worker-LLM calls per completed evaluation. \emph{Information debt} is state debt, lineage debt, stale-winner validation rate, and axis entropy as defined in Section~\ref{sec:problem}.

Best score alone is insufficient for this setting. A method that finds a lucky final improvement after running many stale evaluations may look good by best score but be hard to trust; conversely, a method with slightly worse final best score but lower information debt may be easier to scale.

\subsection{Pilot Autoresearch Results}

Table~\ref{tab:auto} shows the current 1-hour pilot. These numbers should be treated as preliminary because seed-noise estimation and repeated trials are still pending. The qualitative pattern is nevertheless useful: naive parallelism produces many evaluations but weak improvement, while \method recovers a larger improvement by allowing successful edits to enter the shared lineage.

\begin{table}[t]
\centering
\small
\begin{tabular}{lcccc}
\toprule
Policy & Evals & Best \lowerbetter & $\Delta$ & Keeps \\
\midrule
Serial & 8 & 0.9894 & $-0.0015$ & 4 \\
Naive parallel\textsuperscript{$\dagger$} & 27 & 0.9928 & $-0.0001$ & -- \\
\method & 65 & 0.9920 & $-0.0010$ & 1 \\
\bottomrule
\end{tabular}
\caption{Pilot autoresearch results at a 1-hour wall-clock budget. Lower \texttt{val\_bpb} is better. $\Delta$ is improvement over each run's compile-on starting point. Complete repetitions and noise estimates are pending. \textsuperscript{$\dagger$}\,The naive-parallel measurement is from a configuration that also exhibited an unrelated sandbox-isolation bug (later fixed): all parallel evaluations read the same shared source rather than per-evaluation clones. While this exaggerated the metric collapse, even a clean naive-parallel run without lineage cannot produce cumulative wins by construction; we retain the data point as conservative evidence and plan a fixed re-run for camera-ready.}
\label{tab:auto}
\end{table}

The table also exposes a weakness of the current implementation. \method achieved high throughput, but the manager issued only one \keep directive and did so late in the hour (at $t \approx 55$ minutes, with one of eight workers' results still in flight). This means most of the parallel budget still ran on a single baseline, and the kept lineage had no follow-up window in which to compound. The proposed statistical commit gate (Equation~\ref{eq:commit-gate}) and automatic rebase validation are designed to make beneficial keeps happen earlier without overfitting to noise.

\subsection{Planned SFT and Ablation Tables}

The Qwen SFT data-mixture task is intended to test whether \method helps when the search space is lower-dimensional but noisier and where candidate edits are naturally compositional. Table~\ref{tab:sft} is the planned reporting format.

\begin{table}[t]
\centering
\small
\begin{tabular}{lcccc}
\toprule
Policy & Evals & Best val loss \lowerbetter & AUC \higherbetter & Calls/eval \\
\midrule
Serial & \tbd & \tbd & \tbd & \tbd \\
Naive parallel & \tbd & \tbd & \tbd & \tbd \\
\method & \tbd & \tbd & \tbd & \tbd \\
\method + coverage & \tbd & \tbd & \tbd & \tbd \\
\bottomrule
\end{tabular}
\caption{Planned Qwen SFT data-mixture results. The key comparison is whether lineage and coverage steering improve balanced validation loss at the same wall-clock budget.}
\label{tab:sft}
\end{table}

\begin{table}[t]
\centering
\small
\begin{tabular}{lccc}
\toprule
Configuration & Best \lowerbetter & Axis entropy & Manager calls \\
\midrule
Full \method & 0.9920 & \tbd & 14 \\
-- cumulative lineage & 0.9928 & \tbd & \tbd \\
-- rationale ledger & \tbd & \tbd & \tbd \\
-- source injection & \tbd & \tbd & \tbd \\
$\alpha=1.5$ watermark & \tbd & \tbd & \tbd \\
+ coverage manifest & \tbd & \tbd & \tbd \\
\bottomrule
\end{tabular}
\caption{Planned per-channel ablation on autoresearch. The first two rows already suggest that removing lineage collapses to naive-parallel behavior; remaining rows should quantify memory and scheduling costs.}
\label{tab:ablation}
\end{table}

\section{Related Work}
\label{sec:related}

\paragraph{LLMs as optimizers and agents.}
OPRO frames optimization as prompting an LLM with previous solutions and scores \citep{yang2023opro}; PromptAgent uses strategic planning and tree search for prompt optimization \citep{promptagent}; TextGrad backpropagates textual feedback through compound AI systems \citep{textgrad}; and DSPy provides abstractions for compiling and optimizing LM pipelines \citep{dspy}. These works focus primarily on the optimizer or prompt/program representation. \method focuses on the runtime question that arises when such optimizers drive expensive parallel evaluations: what information must be shared across workers and when?

ReAct and LATS study agents that interleave reasoning, action, feedback, and planning \citep{react,lats}. Their environments often expose explicit states or tree nodes. In contrast, agent-led ML experimentation has side effects in a source repository, noisy harnesses, and expensive evaluations. Our git-backed lineage is a practical state representation for this setting.

\paragraph{Autonomous research systems.}
Agent Laboratory and related systems automate larger pieces of the research workflow, including literature review, coding, experimentation, and report writing \citep{agentlab}. \texttt{autoresearch} is narrower and therefore a useful testbed: it isolates the inner loop of code/config proposal, fixed-duration training, metric reading, and keep/revert decisions \citep{karpathy_autoresearch}. \method targets this inner loop and can be used as an execution layer inside broader autonomous research systems.

\paragraph{Parallel hyperparameter optimization.}
Bayesian optimization, Hyperband, ASHA, and population-based training address parallel resource allocation for black-box or low-dimensional hyperparameter spaces \citep{hyperband,asha,pbt}. These methods provide strong baselines when candidates are vectors and the proposal distribution is programmatic. Agent-led search differs because candidates may be structured code edits with natural-language rationales, and a win may alter the source state on which future candidates are meaningful. \method borrows the resource-awareness of these systems but keeps the LLM as a high-level proposer.

\paragraph{Data-mixture optimization.}
Open post-training recipes such as T\"ulu make data sources, evaluation splits, and post-training stages explicit \citep{tulu3}. Automated data-mixture tuning is an attractive domain for agent-led search because mixtures are interpretable, compositional, and often require qualitative reasoning about domains. The planned Qwen SFT experiment tests whether \method helps the agent exploit that structure without collapsing onto one data source.

\section{Discussion}
\label{sec:discussion}

\paragraph{The current bottleneck is not throughput but timely committing.}
Our pilot \method run produced 65 evaluations in one hour, but only one lineage update. This is a poor use of the core mechanism: once a good edit is discovered, every subsequent experiment should usually be launched from a baseline that includes it or from an intentionally different branch. The next implementation should make \keep decisions easier for the manager by presenting a noise-aware threshold and should automatically re-evaluate stale winners on current \HEAD.

\paragraph{Coverage should be programmatic, not only prompted.}
The manager can be asked to diversify, but prompts alone are brittle. Because \texttt{param\_patch} edits are structured, the runtime can compute axis/value coverage exactly. A simple rule such as ``at most half of the next wave may touch the current dominant axis unless that axis has a validated keep'' would directly reduce fixation. We empirically observed the failure mode in our pilot: 40 of 65 \method evaluations were variants of a single \texttt{aspect\_ratio} winner, while several other axes (weight decay, Adam $\beta_2$, scalar LR) were never touched. Coverage steering is especially important for data-mixture search, where nearby mixtures can consume many evaluations without learning much about unexplored domains.

\paragraph{Noise estimation is essential.}
Short training evaluations can differ by seed, compile state, and hardware variance. Without an estimate of $\hat\sigma$, both serial and parallel agents may keep lucky regressions or discard real but small wins. The final workshop version should include repeated baseline evaluations and at least a few paired repeats of kept candidates. This will make the commit gate in Equation~\ref{eq:commit-gate} empirically meaningful.

\paragraph{Failure logs are part of the result.}
Agent-led search produces a rich trace: rationales, failed edits, stale evaluations, manager wake-ups, and lineage decisions. We believe these traces are useful scientific objects. Reporting information-debt metrics alongside best score makes negative results interpretable: a run can fail because the manager proposed bad hypotheses, because the runtime did not compound wins, because the manager fixated on one axis, or because the metric was too noisy.

\section{Limitations}
\label{sec:limitations}

The current evidence is a pilot rather than a complete evaluation. The autoresearch table has not yet been repeated across seeds; the Qwen SFT task is pending; and only the lineage ablation has a filled result. The method also assumes that edits can usually be represented or summarized as structured patches. Free-form code edits are supported, but coverage extraction and rebase validation become harder when edits restructure large portions of the repository. Finally, \method does not solve exploration by itself. It supplies memory and lineage, but the manager may still choose weak hypotheses; hybridizing the LLM manager with bandit or Bayesian proposal filters is a natural extension.

\section{Conclusion}
\label{sec:conclusion}

Serial LLM-agent search is information-rich but compute-poor; naive parallel search is compute-rich but information-poor. \method reframes this as a resource allocation problem over three information channels: state, lineage, and coverage. Watermark-gated replanning controls how quickly new results reach the manager; cumulative lineage lets validated edits compound; and the rationale/coverage ledger preserves the experimental notebook that serial agents implicitly maintain. Pilot autoresearch results indicate that repairing lineage can matter as much as increasing throughput. The next step is to complete noise-aware commits, programmatic coverage steering, and the planned Qwen SFT data-mixture evaluation.

\bibliographystyle{icml2026}
\bibliography{references}

\appendix
\section{Implementation Notes}
\label{app:impl}

Each worker runs in an isolated sandbox cloned from the shared source repository. The wrapper receives an explicit \texttt{REPO\_ROOT}; this prevents accidental evaluation of the shared source instead of the sandbox. Worker processes are launched in detached process groups so timeout cleanup can terminate descendants. A SQLite database with write-ahead logging tracks GPU tokens and task state; JSON event traces store manager rationale, structured edits, metrics, and stderr snippets. The lineage log records \texttt{experiment\_id}, parent commit, resulting commit, manager reason, and timestamp.

The \method implementation is approximately 5{,}000 lines of TypeScript/JavaScript across nine modules (scheduler, watermark gate, lineage manager, ledger, structured-edit applier, sandbox cloner, walltime gate, GPU drain probe, and result validator), plus a Python helper for AST-normalized comparison of constant edits. A unit test suite of 102 cases covers the schema, lineage, edit kinds, and timeout machinery.

\section{Example Task Graph}
\label{app:taskgraph}

\begin{lstlisting}
<!-- TASK_GRAPH -->
{"tasks": [{
  "id": "exp_042_aspect80",
  "execution_mode": "param_patch",
  "base_ref": "HEAD",
  "rationale": "Tests whether a slightly wider FFN ratio improves bpb without increasing depth.",
  "edits": [{"file": "train.py",
             "kind": "constant_replace",
             "name": "ASPECT_RATIO",
             "expected_old_repr": "64",
             "new_repr": "80"}],
  "resources": {"gpus": 1, "cpus": 1},
  "estimated_runtime_seconds": 320,
  "produces_tags": ["metrics:exp_042"]
}]}
<!-- /TASK_GRAPH -->
\end{lstlisting}

\end{document}
